% =============================================================================
% elementos/diagramas/dfd_pipeline.tex
%   Diagrama de Fluxo de Dados arquitetural do pipeline (Cap. 3).
%   Cascata Ingestão -> Detecção (MD) -> Anonimização LGPD -> Classificação
%   (SpeciesNet) -> Re-ID (PetFace/CatFLW). Inclui ramos de descarte
%   (negativos e fauna não-alvo) e o catálogo de indivíduos como reservatório.
%
%   Pacotes assumidos (já carregados em preambulo/pacotes.tex):
%     tikz, shapes.geometric, arrows.meta, positioning, fit, calc,
%     backgrounds, decorations.pathreplacing, shapes.multipart.
%
%   Uso: % =============================================================================
% elementos/diagramas/dfd_pipeline.tex
%   Diagrama de Fluxo de Dados arquitetural do pipeline (Cap. 3).
%   Cascata Ingestão -> Detecção (MD) -> Anonimização LGPD -> Classificação
%   (SpeciesNet) -> Re-ID (PetFace/CatFLW). Inclui ramos de descarte
%   (negativos e fauna não-alvo) e o catálogo de indivíduos como reservatório.
%
%   Pacotes assumidos (já carregados em preambulo/pacotes.tex):
%     tikz, shapes.geometric, arrows.meta, positioning, fit, calc,
%     backgrounds, decorations.pathreplacing, shapes.multipart.
%
%   Uso: % =============================================================================
% elementos/diagramas/dfd_pipeline.tex
%   Diagrama de Fluxo de Dados arquitetural do pipeline (Cap. 3).
%   Cascata Ingestão -> Detecção (MD) -> Anonimização LGPD -> Classificação
%   (SpeciesNet) -> Re-ID (PetFace/CatFLW). Inclui ramos de descarte
%   (negativos e fauna não-alvo) e o catálogo de indivíduos como reservatório.
%
%   Pacotes assumidos (já carregados em preambulo/pacotes.tex):
%     tikz, shapes.geometric, arrows.meta, positioning, fit, calc,
%     backgrounds, decorations.pathreplacing, shapes.multipart.
%
%   Uso: % =============================================================================
% elementos/diagramas/dfd_pipeline.tex
%   Diagrama de Fluxo de Dados arquitetural do pipeline (Cap. 3).
%   Cascata Ingestão -> Detecção (MD) -> Anonimização LGPD -> Classificação
%   (SpeciesNet) -> Re-ID (PetFace/CatFLW). Inclui ramos de descarte
%   (negativos e fauna não-alvo) e o catálogo de indivíduos como reservatório.
%
%   Pacotes assumidos (já carregados em preambulo/pacotes.tex):
%     tikz, shapes.geometric, arrows.meta, positioning, fit, calc,
%     backgrounds, decorations.pathreplacing, shapes.multipart.
%
%   Uso: \input{elementos/diagramas/dfd_pipeline} dentro de uma figure.
% =============================================================================
\begin{figure}[!htbp]
  \centering
  \begin{tikzpicture}[
    >=Stealth,
    node distance=10mm and 12mm,
    every node/.style={font=\footnotesize, align=center},
    proc/.style={rectangle, rounded corners=2pt, draw, thick,
                 minimum width=28mm, minimum height=10mm, fill=blue!5},
    decision/.style={diamond, draw, thick, aspect=2,
                     minimum width=22mm, fill=orange!8, inner sep=1pt},
    store/.style={rectangle split, rectangle split parts=2,
                  rectangle split horizontal=false, draw, thick,
                  minimum width=22mm, fill=gray!10, font=\footnotesize},
    discard/.style={rectangle, draw, thick, dashed,
                    minimum width=22mm, minimum height=8mm, fill=red!5},
    flow/.style={->, thick},
    flowmeta/.style={->, thick, dashed, gray}
  ]
    % Linha 1: ingestão
    \node[proc] (cam) {Câmeras em campo\\(IPCAM/Trail/SBC)};
    \node[proc, right=of cam] (ing) {Fase~I\\Ingestão e\\pré-tratamento};
    \node[store, right=of ing] (raw) {\textbf{Acervo bruto}\nodepart{two}quadros canônicos};


    % Linha 2: detecção + anonimização
    \node[proc, below=of ing] (md) {Fase~II\\MegaDetector\\(detecção)};
    \node[decision, right=of md] (hasperson) {Classe\\``pessoa''?};
    \node[proc, right=of hasperson] (blur) {Mascaramento\\(blur LGPD)};


    % Linha 3: classificação
    \node[proc, below=of md] (sn) {Fase~III\\SpeciesNet\\(classificação)};
    \node[decision, right=of sn] (isfelis) {É \textit{Felis}\\\textit{catus}?};
    \node[discard, right=of isfelis] (nontarget) {Descarte\\fauna não-alvo};


    % Linha 4: re-ID
    \node[proc, below=of sn] (reid) {Fase~IV\\Re-ID\\(PetFace + CatFLW)};
    \node[store, right=of reid] (cat) {\textbf{Catálogo}\nodepart{two}indivíduos da colônia};
    \node[proc, below=of reid] (out) {Relatórios\\AEX Gatos do C2};


    % Descartes laterais (à esquerda)
    \node[discard, left=of md] (empty) {Descarte\\quadros vazios};


    % Setas linha 1
    \draw[flow] (cam) -- (ing);
    \draw[flow] (ing) -- (raw);


    % Seta da store raw para o MD: desce 5mm para o corredor, vai para esquerda, entra no topo do MD
    \draw[flow] (raw.south) -- ++(0,-0.5) -| (md.north);


    % Setas linha 2
    \draw[flow] (md) -- (hasperson);
    \draw[flow] (hasperson) -- node[above,font=\scriptsize]{sim} (blur);


    % Descarte vazios
    \draw[flow] (md.west) -- node[above,font=\scriptsize]{sem animal} (empty.east);


    % Setas da linha 2 para a linha 3 (evitando sobreposição nos blocos isfelis e nontarget)
    % A seta do blur desce 6mm pelo corredor e entra mais à direita do topo do SpeciesNet
    \draw[flow] (blur.south) -- ++(0,-0.6) -| ([xshift=10mm]sn.north);

    % A seta do hasperson desce 3mm (para não conflitar com a linha do blur) e entra no SpeciesNet
    \draw[flow] (hasperson.south) -- ++(0,-0.3) coordinate (hp_mid) -| ([xshift=4mm]sn.north);
    \path (hp_mid) -- ([xshift=4mm]hp_mid -| sn.north) node[midway, above, font=\scriptsize] {não};


    % Setas linha 3
    \draw[flow] (sn) -- (isfelis);
    \draw[flow] (isfelis) -- node[above,font=\scriptsize]{não} (nontarget);

    % Seta isfelis -> reid: desce 5mm pelo corredor e vai para o topo do reid
    \draw[flow] (isfelis.south) -- ++(0,-0.5) coordinate (is_mid) -| (reid.north);
    \path (is_mid) -- (is_mid -| reid.north) node[midway, above, font=\scriptsize] {sim};


    % Setas linha 4 e 5
    \draw[flow] (reid) -- (cat);
    \draw[flowmeta] (cat.south) |- (out.east);
    \draw[flow] (reid) -- (out);


  \end{tikzpicture}
  \caption[Diagrama de fluxo de dados do pipeline de visão computacional]{Diagrama de fluxo de dados (DFD) arquitetural do pipeline de visão computacional. As quatro fases operam em cascata sobre o acervo bruto consolidado pelo ciclo \textit{offline-first}: o MegaDetector descarta quadros vazios e dispara o mascaramento da classe ``pessoa'' antes de qualquer persistência identificável (conformidade com a LGPD, Seção~\ref{sec:pipeline-fase2}); o SpeciesNet descarta fauna não-alvo; e o estágio de reidentificação (re-ID) opera apenas sobre quadros confirmados como \textit{F.\,catus}, contra o catálogo de indivíduos da colônia, alimentando os relatórios operacionais da AEX Gatos do C2.}
  \fonte{Elaborado pelo autor.}
  \label{fig:dfd-pipeline}
\end{figure} dentro de uma figure.
% =============================================================================
\begin{figure}[!htbp]
  \centering
  \begin{tikzpicture}[
    >=Stealth,
    node distance=10mm and 12mm,
    every node/.style={font=\footnotesize, align=center},
    proc/.style={rectangle, rounded corners=2pt, draw, thick,
                 minimum width=28mm, minimum height=10mm, fill=blue!5},
    decision/.style={diamond, draw, thick, aspect=2,
                     minimum width=22mm, fill=orange!8, inner sep=1pt},
    store/.style={rectangle split, rectangle split parts=2,
                  rectangle split horizontal=false, draw, thick,
                  minimum width=22mm, fill=gray!10, font=\footnotesize},
    discard/.style={rectangle, draw, thick, dashed,
                    minimum width=22mm, minimum height=8mm, fill=red!5},
    flow/.style={->, thick},
    flowmeta/.style={->, thick, dashed, gray}
  ]
    % Linha 1: ingestão
    \node[proc] (cam) {Câmeras em campo\\(IPCAM/Trail/SBC)};
    \node[proc, right=of cam] (ing) {Fase~I\\Ingestão e\\pré-tratamento};
    \node[store, right=of ing] (raw) {\textbf{Acervo bruto}\nodepart{two}quadros canônicos};


    % Linha 2: detecção + anonimização
    \node[proc, below=of ing] (md) {Fase~II\\MegaDetector\\(detecção)};
    \node[decision, right=of md] (hasperson) {Classe\\``pessoa''?};
    \node[proc, right=of hasperson] (blur) {Mascaramento\\(blur LGPD)};


    % Linha 3: classificação
    \node[proc, below=of md] (sn) {Fase~III\\SpeciesNet\\(classificação)};
    \node[decision, right=of sn] (isfelis) {É \textit{Felis}\\\textit{catus}?};
    \node[discard, right=of isfelis] (nontarget) {Descarte\\fauna não-alvo};


    % Linha 4: re-ID
    \node[proc, below=of sn] (reid) {Fase~IV\\Re-ID\\(PetFace + CatFLW)};
    \node[store, right=of reid] (cat) {\textbf{Catálogo}\nodepart{two}indivíduos da colônia};
    \node[proc, below=of reid] (out) {Relatórios\\AEX Gatos do C2};


    % Descartes laterais (à esquerda)
    \node[discard, left=of md] (empty) {Descarte\\quadros vazios};


    % Setas linha 1
    \draw[flow] (cam) -- (ing);
    \draw[flow] (ing) -- (raw);


    % Seta da store raw para o MD: desce 5mm para o corredor, vai para esquerda, entra no topo do MD
    \draw[flow] (raw.south) -- ++(0,-0.5) -| (md.north);


    % Setas linha 2
    \draw[flow] (md) -- (hasperson);
    \draw[flow] (hasperson) -- node[above,font=\scriptsize]{sim} (blur);


    % Descarte vazios
    \draw[flow] (md.west) -- node[above,font=\scriptsize]{sem animal} (empty.east);


    % Setas da linha 2 para a linha 3 (evitando sobreposição nos blocos isfelis e nontarget)
    % A seta do blur desce 6mm pelo corredor e entra mais à direita do topo do SpeciesNet
    \draw[flow] (blur.south) -- ++(0,-0.6) -| ([xshift=10mm]sn.north);

    % A seta do hasperson desce 3mm (para não conflitar com a linha do blur) e entra no SpeciesNet
    \draw[flow] (hasperson.south) -- ++(0,-0.3) coordinate (hp_mid) -| ([xshift=4mm]sn.north);
    \path (hp_mid) -- ([xshift=4mm]hp_mid -| sn.north) node[midway, above, font=\scriptsize] {não};


    % Setas linha 3
    \draw[flow] (sn) -- (isfelis);
    \draw[flow] (isfelis) -- node[above,font=\scriptsize]{não} (nontarget);

    % Seta isfelis -> reid: desce 5mm pelo corredor e vai para o topo do reid
    \draw[flow] (isfelis.south) -- ++(0,-0.5) coordinate (is_mid) -| (reid.north);
    \path (is_mid) -- (is_mid -| reid.north) node[midway, above, font=\scriptsize] {sim};


    % Setas linha 4 e 5
    \draw[flow] (reid) -- (cat);
    \draw[flowmeta] (cat.south) |- (out.east);
    \draw[flow] (reid) -- (out);


  \end{tikzpicture}
  \caption[Diagrama de fluxo de dados do pipeline de visão computacional]{Diagrama de fluxo de dados (DFD) arquitetural do pipeline de visão computacional. As quatro fases operam em cascata sobre o acervo bruto consolidado pelo ciclo \textit{offline-first}: o MegaDetector descarta quadros vazios e dispara o mascaramento da classe ``pessoa'' antes de qualquer persistência identificável (conformidade com a LGPD, Seção~\ref{sec:pipeline-fase2}); o SpeciesNet descarta fauna não-alvo; e o estágio de reidentificação (re-ID) opera apenas sobre quadros confirmados como \textit{F.\,catus}, contra o catálogo de indivíduos da colônia, alimentando os relatórios operacionais da AEX Gatos do C2.}
  \fonte{Elaborado pelo autor.}
  \label{fig:dfd-pipeline}
\end{figure} dentro de uma figure.
% =============================================================================
\begin{figure}[!htbp]
  \centering
  \begin{tikzpicture}[
    >=Stealth,
    node distance=10mm and 12mm,
    every node/.style={font=\footnotesize, align=center},
    proc/.style={rectangle, rounded corners=2pt, draw, thick,
                 minimum width=28mm, minimum height=10mm, fill=blue!5},
    decision/.style={diamond, draw, thick, aspect=2,
                     minimum width=22mm, fill=orange!8, inner sep=1pt},
    store/.style={rectangle split, rectangle split parts=2,
                  rectangle split horizontal=false, draw, thick,
                  minimum width=22mm, fill=gray!10, font=\footnotesize},
    discard/.style={rectangle, draw, thick, dashed,
                    minimum width=22mm, minimum height=8mm, fill=red!5},
    flow/.style={->, thick},
    flowmeta/.style={->, thick, dashed, gray}
  ]
    % Linha 1: ingestão
    \node[proc] (cam) {Câmeras em campo\\(IPCAM/Trail/SBC)};
    \node[proc, right=of cam] (ing) {Fase~I\\Ingestão e\\pré-tratamento};
    \node[store, right=of ing] (raw) {\textbf{Acervo bruto}\nodepart{two}quadros canônicos};


    % Linha 2: detecção + anonimização
    \node[proc, below=of ing] (md) {Fase~II\\MegaDetector\\(detecção)};
    \node[decision, right=of md] (hasperson) {Classe\\``pessoa''?};
    \node[proc, right=of hasperson] (blur) {Mascaramento\\(blur LGPD)};


    % Linha 3: classificação
    \node[proc, below=of md] (sn) {Fase~III\\SpeciesNet\\(classificação)};
    \node[decision, right=of sn] (isfelis) {É \textit{Felis}\\\textit{catus}?};
    \node[discard, right=of isfelis] (nontarget) {Descarte\\fauna não-alvo};


    % Linha 4: re-ID
    \node[proc, below=of sn] (reid) {Fase~IV\\Re-ID\\(PetFace + CatFLW)};
    \node[store, right=of reid] (cat) {\textbf{Catálogo}\nodepart{two}indivíduos da colônia};
    \node[proc, below=of reid] (out) {Relatórios\\AEX Gatos do C2};


    % Descartes laterais (à esquerda)
    \node[discard, left=of md] (empty) {Descarte\\quadros vazios};


    % Setas linha 1
    \draw[flow] (cam) -- (ing);
    \draw[flow] (ing) -- (raw);


    % Seta da store raw para o MD: desce 5mm para o corredor, vai para esquerda, entra no topo do MD
    \draw[flow] (raw.south) -- ++(0,-0.5) -| (md.north);


    % Setas linha 2
    \draw[flow] (md) -- (hasperson);
    \draw[flow] (hasperson) -- node[above,font=\scriptsize]{sim} (blur);


    % Descarte vazios
    \draw[flow] (md.west) -- node[above,font=\scriptsize]{sem animal} (empty.east);


    % Setas da linha 2 para a linha 3 (evitando sobreposição nos blocos isfelis e nontarget)
    % A seta do blur desce 6mm pelo corredor e entra mais à direita do topo do SpeciesNet
    \draw[flow] (blur.south) -- ++(0,-0.6) -| ([xshift=10mm]sn.north);

    % A seta do hasperson desce 3mm (para não conflitar com a linha do blur) e entra no SpeciesNet
    \draw[flow] (hasperson.south) -- ++(0,-0.3) coordinate (hp_mid) -| ([xshift=4mm]sn.north);
    \path (hp_mid) -- ([xshift=4mm]hp_mid -| sn.north) node[midway, above, font=\scriptsize] {não};


    % Setas linha 3
    \draw[flow] (sn) -- (isfelis);
    \draw[flow] (isfelis) -- node[above,font=\scriptsize]{não} (nontarget);

    % Seta isfelis -> reid: desce 5mm pelo corredor e vai para o topo do reid
    \draw[flow] (isfelis.south) -- ++(0,-0.5) coordinate (is_mid) -| (reid.north);
    \path (is_mid) -- (is_mid -| reid.north) node[midway, above, font=\scriptsize] {sim};


    % Setas linha 4 e 5
    \draw[flow] (reid) -- (cat);
    \draw[flowmeta] (cat.south) |- (out.east);
    \draw[flow] (reid) -- (out);


  \end{tikzpicture}
  \caption[Diagrama de fluxo de dados do pipeline de visão computacional]{Diagrama de fluxo de dados (DFD) arquitetural do pipeline de visão computacional. As quatro fases operam em cascata sobre o acervo bruto consolidado pelo ciclo \textit{offline-first}: o MegaDetector descarta quadros vazios e dispara o mascaramento da classe ``pessoa'' antes de qualquer persistência identificável (conformidade com a LGPD, Seção~\ref{sec:pipeline-fase2}); o SpeciesNet descarta fauna não-alvo; e o estágio de reidentificação (re-ID) opera apenas sobre quadros confirmados como \textit{F.\,catus}, contra o catálogo de indivíduos da colônia, alimentando os relatórios operacionais da AEX Gatos do C2.}
  \fonte{Elaborado pelo autor.}
  \label{fig:dfd-pipeline}
\end{figure} dentro de uma figure.
% =============================================================================
\begin{figure}[!htbp]
  \centering
  \begin{tikzpicture}[
    >=Stealth,
    node distance=10mm and 12mm,
    every node/.style={font=\footnotesize, align=center},
    proc/.style={rectangle, rounded corners=2pt, draw, thick,
                 minimum width=28mm, minimum height=10mm, fill=blue!5},
    decision/.style={diamond, draw, thick, aspect=2,
                     minimum width=22mm, fill=orange!8, inner sep=1pt},
    store/.style={rectangle split, rectangle split parts=2,
                  rectangle split horizontal=false, draw, thick,
                  minimum width=22mm, fill=gray!10, font=\footnotesize},
    discard/.style={rectangle, draw, thick, dashed,
                    minimum width=22mm, minimum height=8mm, fill=red!5},
    flow/.style={->, thick},
    flowmeta/.style={->, thick, dashed, gray}
  ]
    % Linha 1: ingestão
    \node[proc] (cam) {Câmeras em campo\\(IPCAM/Trail/SBC)};
    \node[proc, right=of cam] (ing) {Fase~I\\Ingestão e\\pré-tratamento};
    \node[store, right=of ing] (raw) {\textbf{Acervo bruto}\nodepart{two}quadros canônicos};


    % Linha 2: detecção + anonimização
    \node[proc, below=of ing] (md) {Fase~II\\MegaDetector\\(detecção)};
    \node[decision, right=of md] (hasperson) {Classe\\``pessoa''?};
    \node[proc, right=of hasperson] (blur) {Mascaramento\\(blur LGPD)};


    % Linha 3: classificação
    \node[proc, below=of md] (sn) {Fase~III\\SpeciesNet\\(classificação)};
    \node[decision, right=of sn] (isfelis) {É \textit{Felis}\\\textit{catus}?};
    \node[discard, right=of isfelis] (nontarget) {Descarte\\fauna não-alvo};


    % Linha 4: re-ID
    \node[proc, below=of sn] (reid) {Fase~IV\\Re-ID\\(PetFace + CatFLW)};
    \node[store, right=of reid] (cat) {\textbf{Catálogo}\nodepart{two}indivíduos da colônia};
    \node[proc, below=of reid] (out) {Relatórios\\AEX Gatos do C2};


    % Descartes laterais (à esquerda)
    \node[discard, left=of md] (empty) {Descarte\\quadros vazios};


    % Setas linha 1
    \draw[flow] (cam) -- (ing);
    \draw[flow] (ing) -- (raw);


    % Seta da store raw para o MD: desce 5mm para o corredor, vai para esquerda, entra no topo do MD
    \draw[flow] (raw.south) -- ++(0,-0.5) -| (md.north);


    % Setas linha 2
    \draw[flow] (md) -- (hasperson);
    \draw[flow] (hasperson) -- node[above,font=\scriptsize]{sim} (blur);


    % Descarte vazios
    \draw[flow] (md.west) -- node[above,font=\scriptsize]{sem animal} (empty.east);


    % Setas da linha 2 para a linha 3 (evitando sobreposição nos blocos isfelis e nontarget)
    % A seta do blur desce 6mm pelo corredor e entra mais à direita do topo do SpeciesNet
    \draw[flow] (blur.south) -- ++(0,-0.6) -| ([xshift=10mm]sn.north);

    % A seta do hasperson desce 3mm (para não conflitar com a linha do blur) e entra no SpeciesNet
    \draw[flow] (hasperson.south) -- ++(0,-0.3) coordinate (hp_mid) -| ([xshift=4mm]sn.north);
    \path (hp_mid) -- ([xshift=4mm]hp_mid -| sn.north) node[midway, above, font=\scriptsize] {não};


    % Setas linha 3
    \draw[flow] (sn) -- (isfelis);
    \draw[flow] (isfelis) -- node[above,font=\scriptsize]{não} (nontarget);

    % Seta isfelis -> reid: desce 5mm pelo corredor e vai para o topo do reid
    \draw[flow] (isfelis.south) -- ++(0,-0.5) coordinate (is_mid) -| (reid.north);
    \path (is_mid) -- (is_mid -| reid.north) node[midway, above, font=\scriptsize] {sim};


    % Setas linha 4 e 5
    \draw[flow] (reid) -- (cat);
    \draw[flowmeta] (cat.south) |- (out.east);
    \draw[flow] (reid) -- (out);


  \end{tikzpicture}
  \caption[Diagrama de fluxo de dados do pipeline de visão computacional]{Diagrama de fluxo de dados (DFD) arquitetural do pipeline de visão computacional. As quatro fases operam em cascata sobre o acervo bruto consolidado pelo ciclo \textit{offline-first}: o MegaDetector descarta quadros vazios e dispara o mascaramento da classe ``pessoa'' antes de qualquer persistência identificável (conformidade com a LGPD, Seção~\ref{sec:pipeline-fase2}); o SpeciesNet descarta fauna não-alvo; e o estágio de reidentificação (re-ID) opera apenas sobre quadros confirmados como \textit{F.\,catus}, contra o catálogo de indivíduos da colônia, alimentando os relatórios operacionais da AEX Gatos do C2.}
  \fonte{Elaborado pelo autor.}
  \label{fig:dfd-pipeline}
\end{figure}